%%=============================================================================
%% System Usability Scale (SUS) - Nederlandstalige vragenlijst
%%
%% Standaard 10-item SUS volgens Brooke (1996). Schaal: 1 (helemaal mee oneens)
%% tot 5 (helemaal mee eens). Items 1, 3, 5, 7, 9 zijn positief geformuleerd.
%% items 2, 4, 6, 8, 10 zijn negatief geformuleerd.
%%
%% Scoring (na afname):
%%   - Voor positieve items (oneven): score - 1
%%   - Voor negatieve items (even):   5 - score
%%   - Som * 2,5 = SUS-score (0-100)
%%
%% Compileren: xelatex sus-vragenlijst.tex
%%=============================================================================

\documentclass[11pt,a4paper]{article}

\usepackage[dutch]{babel}
\usepackage[utf8]{inputenc}
\usepackage[T1]{fontenc}
\usepackage[margin=2.5cm]{geometry}
\usepackage{xcolor}
\usepackage{enumitem}
\usepackage{titlesec}
\usepackage{fancyhdr}
\usepackage{tabularx}
\usepackage{array}
\usepackage{amssymb}
\usepackage{hyperref}

\definecolor{cobolblue}{RGB}{0,80,140}

\pagestyle{fancy}
\fancyhf{}
\fancyhead[L]{\textsc{Learn COBOL \textbar\ SUS-vragenlijst}}
\fancyhead[R]{Pagina \thepage}
\fancyfoot[C]{Bachelorproef Silas De Vuyst, HOGENT, 2025--2026}

\titleformat{\section}{\Large\bfseries\color{cobolblue}}{\thesection}{1em}{}

\hypersetup{colorlinks=true,linkcolor=cobolblue,urlcolor=cobolblue}

%% Een handig commando voor één Likert-rij
\newcommand{\likertrow}[2]{%
    \begin{tabularx}{\textwidth}{@{}p{0.7cm} X *{5}{>{\centering\arraybackslash}p{1.1cm}}@{}}
    \textbf{#1.} & #2 & $\square$ 1 & $\square$ 2 & $\square$ 3 & $\square$ 4 & $\square$ 5 \\
    \end{tabularx}\par\medskip%
}

\begin{document}

%%---------- Titelblok -------------------------------------------------------
\begin{center}
{\LARGE\bfseries System Usability Scale (SUS)}\\[0.3cm]
{\large \textit{Vragenlijst voor de gebruikersstudie ``Learn COBOL''}}\\[0.6cm]
\rule{0.6\textwidth}{0.4pt}
\end{center}

\bigskip

\noindent\textbf{Deelnemernummer:}\ \rule{3cm}{0.4pt}\hfill\textbf{Datum:}\ \rule{3cm}{0.4pt}

\bigskip

\noindent
Bedankt om \textit{Learn COBOL} te testen. Onderstaande vragenlijst peilt naar de \textbf{bruikbaarheid} (\textit{usability}) van de applicatie. 
Geef voor elke uitspraak aan in welke mate je het ermee eens bent op een schaal van 1 tot 5:

\begin{center}
\begin{tabularx}{0.85\textwidth}{X X X X X}
\textbf{1} & \textbf{2} & \textbf{3} & \textbf{4} & \textbf{5} \\
\hline
\textit{Helemaal} & \textit{Eerder} & \textit{Neutraal} & \textit{Eerder} & \textit{Helemaal} \\
\textit{mee oneens} & \textit{mee oneens} & & \textit{mee eens} & \textit{mee eens} \\
\end{tabularx}
\end{center}

\bigskip

\noindent
Er zijn geen ``juiste'' of ``foute'' antwoorden. Antwoord op elke uitspraak en denk niet te lang na: je eerste reactie volstaat. Vink één bolletje aan per uitspraak.

%%---------- 10 SUS-items ---------------------------------------------------
\section*{Vragenlijst}

\likertrow{1}{Ik denk dat ik dit systeem vaak zou willen gebruiken.}
\likertrow{2}{Ik vond het systeem onnodig complex.}
\likertrow{3}{Ik vond het systeem eenvoudig te gebruiken.}
\likertrow{4}{Ik denk dat ik de hulp van een technische persoon nodig zou hebben om dit systeem te kunnen gebruiken.}
\likertrow{5}{Ik vond dat de verschillende functies in dit systeem goed waren geïntegreerd.}
\likertrow{6}{Ik vond dat er te veel inconsistentie in dit systeem zat.}
\likertrow{7}{Ik kan me voorstellen dat de meeste mensen het gebruik van dit systeem snel onder de knie zouden krijgen.}
\likertrow{8}{Ik vond het systeem erg omslachtig in gebruik.}
\likertrow{9}{Ik voelde me erg zelfverzekerd bij het gebruik van het systeem.}
\likertrow{10}{Ik moest veel dingen leren voordat ik aan de slag kon met dit systeem.}

\bigskip

\noindent\textbf{Optionele open vraag.} Wat zou jij als allereerste aanpassen aan \textit{Learn COBOL}?

\medskip
\noindent\rule{\textwidth}{0.4pt}\\[0.4cm]
\rule{\textwidth}{0.4pt}\\[0.4cm]
\rule{\textwidth}{0.4pt}

\bigskip

\noindent\textit{Bedankt voor je tijd en je eerlijke antwoord.}

%%---------- Scoring (alleen voor onderzoeker) -------------------------------
\newpage
\section*{Scoring (niet voor de deelnemer)}

\noindent
SUS-score wordt als volgt berekend:

\begin{enumerate}
    \item Voor de \textbf{oneven} items (1, 3, 5, 7, 9: positief geformuleerd): score $-$ 1.
    \item Voor de \textbf{even} items (2, 4, 6, 8, 10: negatief geformuleerd): 5 $-$ score.
    \item Tel de tien aangepaste waarden op (range 0--40).
    \item Vermenigvuldig met 2,5 om een score op 100 te krijgen.
\end{enumerate}

\noindent
Interpretatieschaal volgens \textit{Bangor, Kortum \& Miller (2009)}:

\begin{center}
\begin{tabular}{|c|l|l|}
\hline
\textbf{SUS-score} & \textbf{Adjectief} & \textbf{Letter-grade} \\
\hline
< 50    & Awful / Poor   & F \\
\hline
50--67  & OK             & D \\
\hline
68      & Industriegemiddelde & C \\
\hline
68--80  & Good           & B \\
\hline
> 80    & Excellent      & A \\
\hline
\end{tabular}
\end{center}

\bigskip

\noindent\textbf{Bron.} Brooke, J.\ (1996). \textit{SUS: A quick and dirty usability scale}. In Jordan, P.\ et al.\ (red.), \textit{Usability evaluation in industry}, p.\ 189--194.

\end{document}
